\begin{figure}[H]
\centering
\begin{tikzpicture}[
  >=Latex,
  node distance=7mm,
  box/.style={draw, rounded corners=2pt, align=center, minimum height=30mm,
    text width=30mm, inner sep=4pt, font=\small},
  flow/.style={->, semithick},
]
  \node[box, fill=blue!8] (input) {finite source model\\target and contexts\\supports and resources};
  \node[box, fill=green!9, right=of input] (compile) {admitted state\\labels, spans, classes\\support choices\\symmetry and loads};
  \node[box, fill=orange!10, right=of compile] (solve) {exact engines\\composition and search\\pricing and scheduling\\admitted updates};
  \node[box, fill=violet!9, right=of solve] (output) {answer and evidence\\coefficient witnesses\\bounds and separators\\scoped verification};
  \draw[flow] (input) -- (compile);
  \draw[flow] (compile) -- (solve);
  \draw[flow] (solve) -- (output);
\end{tikzpicture}
\caption{Ergodis separates the source model, admitted representation,
optimization, and evidence contract.  Functional labels preserve composition;
support alternatives preserve overlap and helper-price queries.  Stored lifts
reconstruct feasible decisions, while optimality needs a matching checked
bound or complete search justification.  Reuse under updates or new readouts
requires admission for the chosen observation.  The mathematical results do
not depend on executing this pipeline.}
\label{fig:ergodis-pipeline}
\end{figure}
